\textbf{ML weather forecasting and benchmarks.} WeatherBench and WeatherBench~2 standardise verification of data-driven forecasts with RMSE and ACC against persistence, climatology and physical models~\cite{rasp2020weatherbench,rasp2023weatherbench}; we adopt the per-lead reporting and the two trivial baselines. GraphCast, FourCastNet and the graph network of Keisler are autoregressive emulators at global scale~\cite{lam2022graphcast,pathak2022fourcastnet,keisler2022forecasting}; NeuralGCM couples a learned component to a differentiable dynamical core and is trained through rollouts~\cite{kochkov2023neural}. To our knowledge their rollout-length choices are made at full scale together with other design changes, and none can be re-run here; our CNN is a stand-in that shares only the autoregressive structure.

\textbf{Lorenz-96 as a testbed.} Chattopadhyay et al. compare reservoir computing, feed-forward networks and LSTMs for forecasting the slow variables of a multi-scale Lorenz-96 system~\cite{chattopadhyay2019data}; Gagne et al. learn a stochastic parameterisation of the fast variables with a GAN~\cite{gagne2019machine}. Our emulator is deterministic and acts on the slow variables only, so unresolved-scale effects enter as irreducible one-step error.

\textbf{Stable long rollouts.} Learned coarse turbulence models and neural PDE solvers use training noise and the pushforward trick to counter error accumulation~\cite{stachenfeld2021learned,brandstetter2022message}; PDE-Refiner replaces a single deterministic step by iterative refinement so that low-amplitude, high-frequency modes are not lost~\cite{lippe2023pde}; Thermalizer corrects emulator states back toward the invariant distribution using a diffusion model~\cite{pedersen2025thermalizer}. We test the simplest member of this family, training noise, and a multi-step loss, and report that in our setting the unmodified one-step model needed neither.

\textbf{Gap.} We found no cheap, multi-seed experiment that reports skill (RMSE, ACC versus lead), survival of long free runs and the climatological variance ratio together as a function of $K$; this paper provides one. We do not claim it is the first.
